% answers_short.tex —— 2008 年·普通高等学校招生考试·北京卷·理科 · 简短答案速查
% 由 tools/make_short_answers.py 自 answers.tex 的 \gksolhead/\gkanswer 提取，
% 并与 2008/2008_数学真题_北京理_答案要点.md 逐题比对后生成，勿手改（改源 MD 后重跑）。

\begin{gknotice}
\textbf{使用说明}：本节只给最终答案便于快速核对，完整推导见解析版。答案为编写者独立推导，\textbf{非官方标准答案}。
\end{gknotice}

\gkpart{选择题}
\begin{center}\begin{tabular}{ccccccccc}
\toprule
题号 & 1 & 2 & 3 & 4 & 5 & 6 & 7 & 8 \\
\midrule
答案 & D & A & B & D & B & C & C & B \\
\bottomrule
\end{tabular}\end{center}

\gkpart{填空题}
\begin{itemize}[leftmargin=2.6em]
  \item[\textbf{9.}] $-1$。
  \item[\textbf{10.}] $0$。
  \item[\textbf{11.}] $n=5$；$10$。
  \item[\textbf{12.}] $2$；$-2$。
  \item[\textbf{13.}] ②。
  \item[\textbf{14.}] $(1,2)$；$(3,402)$。
\end{itemize}

\gkpart{解答题}
\begin{itemize}[leftmargin=2.6em]
  \item[\textbf{15.}] （1）$\omega=1$；（2）$\left[0,\dfrac32\right]$。
  \item[\textbf{16.}] （1）证明见解析；（2）$\arccos\dfrac1{\sqrt3}$；（3）$\dfrac{2\sqrt3}3$。
  \item[\textbf{17.}] （1）$\dfrac1{40}$；（2）$\dfrac9{10}$；（3）$P(\xi=1)=\dfrac34$，$P(\xi=2)=\dfrac14$。
  \item[\textbf{18.}] $f'(x)=-\dfrac{2(x-b+1)}{(x-1)^3}$；单调区间按 $b>2$、$b=2$、$b<2$ 分列，见解析。
  \item[\textbf{19.}] （1）$x+y+2=0$；（2）$4\sqrt3$。
  \item[\textbf{20.}] （1）$A_1=A_2=(4,3,2,1)$；（2）（3）证明见解析。
\end{itemize}
